\section{工程化蓝图：把课程观点映射为可执行研发流程}

\subsection{第一步：先写清楚 ``你要对谁最优''}
很多团队在项目立项时直接写 ``提升 Agent 成功率''，但没有定义成功率针对的群体。按照本讲观点，这会在多人任务中造成根本性偏差。一个可执行写法应至少包含以下字段：
\begin{itemize}
    \item 目标群体：同类 Agent、人类用户、还是混合群体；
    \item 目标行为：稳健不亏、平均收益最大，还是合作满意度最大；
    \item 约束条件：时延、成本、可解释性、安全策略上限。
\end{itemize}

\begin{importantbox}{立项阶段的硬要求}
若任务是 ``与人协作''，训练分布里必须显式包含目标人群行为样本；否则上线后的群体错配几乎必然发生。
\end{importantbox}

对应到课程原话，所谓 ``population'' 不是抽象概念，而是必须落在数据治理文档中的实体对象。谁是 population，谁就决定了策略最优性的定义边界。

\subsection{第二步：训练链路按三层解耦}
结合 Brown 在 Diplomacy/Hanabi 的经验，推荐把训练流程拆成三层，每层单独验收：
\begin{enumerate}
    \item \textbf{Behavior Layer}：imitation model，先拟合目标群体可观测行为；
    \item \textbf{Inference Layer}：放大 test-time compute，提高行为预测与策略评估精度；
    \item \textbf{Policy Layer}：在前两层构建的环境分布上做 RL 优化。
\end{enumerate}

\begin{knowledgebox}{为什么要分层验收}
若把 imitation、search、RL 全揉进一个端到端指标，模型退化时很难定位是 ``群体建模失败'' 还是 ``策略优化失败''。分层验收能把排障成本降一个量级。
\end{knowledgebox}

一个可复现实验模板如下：
\begin{lstlisting}[caption={面向多人协作任务的三层训练模板}]
# Layer 1: fit target population
imitation_model = train_imitation(human_or_target_data)

# Layer 2: improve test-time modeling
planner = build_inference_scaler(imitation_model, search_budget=B)

# Layer 3: optimize policy in modeled population
policy = train_rl(env_with(planner), objective="population_utility")
evaluate(policy, cross_population_benchmarks)
\end{lstlisting}

\begin{warningbox}{端到端捷径的代价}
直接在稀疏奖励上做 RL，短期可能拿到局部高分，但会把 ``错误人群假设'' 固化进策略，后期修复成本极高。
\end{warningbox}

\subsection{第三步：评测矩阵必须覆盖跨群体迁移}
课程里 DORA 与 SearchBot 的交叉失败，本质上就是评测矩阵不完整的反例教材。最低建议是三轴评测：
\begin{itemize}
    \item 同群体（in-population）表现；
    \item 异群体（out-of-population）表现；
    \item 长回合稳定性（multi-turn stability）表现。
\end{itemize}

\begin{table}[H]
\centering
\begin{tabular}{p{3.2cm}p{3.0cm}p{3.0cm}p{3.8cm}}
\toprule
\textbf{评测轴} & \textbf{核心指标} & \textbf{报警阈值示例} & \textbf{诊断意义} \\
\midrule
同群体效能 & 平均收益/胜率 & 相比基线提升 $<2\%$ & 训练是否学到任务主干能力 \\
异群体鲁棒性 & 跨群体收益掉点 & 掉点 $>15\%$ & 是否过拟合单一均衡盆地 \\
长程一致性 & 多轮目标一致率 & 低于 85\% & 协商机制是否可持续 \\
资源效率 & token/时延成本 & 成本翻倍收益 $<5\%$ & 并行扩展是否具有性价比 \\
\bottomrule
\end{tabular}
\caption{多 Agent 系统上线前的最低评测矩阵}
\end{table}

\subsection{第四步：上线策略从 ``强模型'' 转向 ``稳协议''}
课程后半段关于 scaffold 脆弱性的讨论，实操上意味着架构重点要从 ``再加一个 agent'' 转为 ``协议正确性''。建议至少具备：
\begin{itemize}
    \item 统一任务语义层（shared task schema）；
    \item 冲突仲裁层（conflict resolver）；
    \item 回退执行层（fallback to single-agent safe mode）。
\end{itemize}

这三层并不华丽，但通常比再加一个复杂协作循环更能提升生产可靠性。

\subsection{本章小结}
课程观点可以落成一条研发流水线：\textbf{群体定义 $\rightarrow$ 分层训练 $\rightarrow$ 跨群体评测 $\rightarrow$ 协议化上线}。这个顺序不应颠倒。

